% Ad Familiares 12.13 (ad-familiares-12-13)
% Source: https://raw.githubusercontent.com/PerseusDL/canonical-latinLit/master/data/phi0474/phi056/phi0474.phi056.perseus-lat1.xml
% Fetched by scripts/fetch_latin.py

% Dateline (Perseus): Scr. in Cypro Crommyuacride Id. Iun. a. 711 (43) .

\ciceroLetterOpener{C. CASSIVS Q. S. D. M. CICERONI}

\ciceroSection{1} S. v. b. e. v. Cum rei p. vel salute vel victoria gaudemus tum instauratione tuarum laudum, quod maximus consularis maximum consulem te ipse vicisti, et laetamur et mirari satis non possumus. fatale nescio quid tuae virtuti datum, id quod saepe iam experti sumus. est enim tua toga omnium armis felicior ; quae nunc quoque nobis paene victam rem p. ex manibus hostium eripuit ac reddidit. nunc ergo vivemus liberi, nunc te, omnium maxime civis et mihi carissime, id quod maximis rei p. tenebris comperisti, nunc te habebimus testem nostri et in te et in coniunctissimam tibi rem p. amoris et, quae saepe pollicitus es te et taciturum dum serviremus, et dicturum de me tum cum mihi profutura essent, nunc illa non ego quidem dici tanto opere desiderabo quam sentiri a te ipso. neque enim omnium iudicio malim me a te commendari quam ipse tuo iudicio digne ac mereor commendatus esse, ut haec novissima nostra facta non subita nec convenientia sed similia illis cogitationibus quarum tu testis es fuisse iudices meque ad optimam spem patriae non minimum tibi ipsi producendum putes.

\ciceroSection{2} sunt tibi, M. Tulli, liberi propinquique digni quidem te et merito tibi carissimi ; esse etiam debent in re p. proxime hos can, qui studiorum tuorum sunt aemuli, quorum esse cupio tibi copiam ; sed tamen non maxima me turba puto excludi, quo minus tibi vacet me excipere et ad omnia quae velis et probes producere. animum tibi nostrum fortasse probavimus, ingenium diutina servitus certe, qualecumque est, minus tamen quam erat passa est videri.

\ciceroSection{3} nos ex ora maritima Asiae provinciae et ex insulis quas potuimus navis deduximus, dilectum remigum magna contumacia civitatium tamen satis celeriter habuimus, secuti sumus classem Dolabellae cui L. Figulus praeerat ; qui spem saepe transitionis praebendo neque umquam non decedendo novissime Corycum se contulit et clauso portu se tenere coepit. nos illa relicta, quod et in castra pervenire satius esse putabamus et sequebatur classis altera, quam anno priore in Bithynia Tillius Cimber compararat, Turullius quaestor praeerat, Cyprum petivimus. ibi quae cognovimus scribere ad vos quam celerrime voluiiuus.

\ciceroSection{4} Dolabellam ut Tarsenses, pessimi socii, ita Laudiceni multo amentiores ultro arcessierunt; ex quibus utrisque civitatibus Graecorum militum numero speciem exercitus effecit. castra habet ante oppidum Laudiceam posita et partem muri demolitus est et castra oppido coniunxit. Cassius noster cum decem legionibus et cohortibus xx auxiliariis et quattuor milium equitatu a milibus passuum xx castra habet posita *pa/ltw| et existimat se sine proelio posse vincere ; nam iam ternis tetrachmis triticum apud Dolabellam est. Nisi quid navibus Laudicenorum supportarit, cito fame pereat necesse est ; ne supportare possit et Cassi classis bene magna cui praeest Sextilius Rufus et tres quas nos adduximus, ego, Turullius, Patiscus, facile praestabunt. te volo bene sperare et rem p., ut vos istic expedistis, ita pro nostra parte celeriter nobis expediri posse confidere. vale . D. Idib. Iun. Cypro a Crommyuacride.
